% =============================================================================
% §7 Scope & limitations (M5) — inferential data-limitations subsection;
% claims-discipline constraints stated as scope conditions.
% =============================================================================
\section{Scope and Limitations}
\label{sec:limitations}

We state the paper's boundaries as claims we do \emph{not} make, each tied to
the data or design feature that enforces it.

\subsection{Data limitations: inference}
\label{sec:datalim-inference}

\paragraph{Pseudo-panels, not transitions.} All state-dependence estimates rest
on repeated cross-sections: within-pseudo-cohort cell dynamics, individual
\emph{outcomes} (Section~\ref{sec:individual}), but never observed individual
\emph{transitions}. First-union timing shifts can masquerade as composition
change in a pseudo-panel; the tempo margin is flagged, not resolved. And the
limitation should be stated at full strength: a no-reflexivity model sufficing
at the annual cell level does not \emph{close} the amplification question ---
it \emph{relocates} it. If a coordination dynamic operates anywhere, it would
live in first-union \emph{timing}, at sub-annual resolution, among the 15--19
entrants our estimation floor never observes --- scales at which annual
repeated cross-sections are blind by construction. The Colombian ENDS
event-history program is therefore not merely the resolution of the 2019--22
shape residual (Section~\ref{sec:results}); it is \emph{the test of sub-annual
amplification that the annual data cannot perform}. That is the honest content
of ``reflexivity not disproven.''

\paragraph{The 15--19 entry margin is unobserved.} The 20--39 estimation floor
excludes teenagers --- the plausible vanguard of any norm cascade. A cascade
igniting at 15--19 and arriving in our window already converted into cohort
heterogeneity would be under-observed. This is the one qualification we place
on the non-cascade verdict, and it is testable with entry-margin data.

\paragraph{Reflexivity: not warranted, not disproven.} Our design tests the
within-cohort, sign-based cascade signature, and the data reject it. We do not
claim to have disproven all reflexive mechanisms --- reflection-problem
identification limits are real \citep{manski1993} --- only that the mechanism's
observable signature is absent where it should appear, at both the cell and
individual-outcome level, in both countries. Transition-level data can reopen
the question; nothing else in this paper should.

\subsection{Scope conditions on the claims}

\paragraph{Composition, not fertility magnitude.} The sufficiency result of
Section~\ref{sec:model} concerns the \emph{union-composition} collapse. We make
no claim that the identified mechanism quantitatively accounts for the TFR
collapse: that requires union-specific fertility schedules (in particular the
cohabiting-to-married fertility differential), for which the required
married-specific rates are still under acquisition. The generated-TFR overlay
shown with the model is comparison-only by construction (the identification
wall of Section~\ref{sec:calibration}).

\paragraph{Two high-cohabitation regimes.} The identification and sufficiency
results cover Costa Rica and Colombia. External validity to marriage-dominant
regimes is precisely the out-of-sample prediction of
Section~\ref{sec:interpretation}, not an assumption --- Argentina and Chile are
the designated tests.

\paragraph{No tempo decomposition.} The model's fertility layer is
parity-independent and makes no tempo claim; where TFR levels are anchored, a
tempo-adjusted sensitivity in the spirit of \citet{bongaarts1998} is reported
alongside (Section~\ref{sec:datalim-measurement}).

\paragraph{Aggregation honesty.} The Colombian individual-level confirmation
covers outcomes, not transitions; Costa Rican estimates remain cell-level
(pre-tabulated REDATAM aggregates). We flag rather than smooth this asymmetry;
the verdict is sign-based and agrees across every level at which it can be
computed.

\paragraph{Interpretation is severable.} The channel-assisted-compression frame
(Section~\ref{sec:interpretation}) organizes the results and generates the
prediction; the measured collapse, the non-cascade verdict, and the sufficiency
result do not depend on it.
